\section{Fibration axioms, modalities, and exactness}

\begin{claim}{E.1}{Known}{The accessible localization attached to a localizer}
For an accessible fundamental localizer $\W$, let
$L:\cS\rightleftarrows\cS_{\W}:\iota$ be the accessible reflective
localization of spaces corresponding to $\W$. A functor $u$ is in
$\W$ if and only if $L(Bu)$ is an equivalence. We identify
$\cS_{\W}$ with its full subcategory of local spaces. Limits of
local spaces are computed in $\cS$; colimits in $\cS_{\W}$ are
obtained by applying $L$ to colimits in $\cS$. In particular $*$
is local. These conventions concern the infinity-category, not
just its homotopy category.
\end{claim}
The passage from fundamental localizers to localizers of simplicial
sets and the preservation of accessibility are
\citet[Theorem 4.2.15, Remark 4.2.16]{RootCisinski2006}. The associated
model structures are left Bousfield localizations of the Kan model
structure. Their derived localizations give the stated reflective
localization of spaces. The unit is denoted
$\eta_X:X\to LX$ for every space $X$. No such reflector is being
assumed for an arbitrary nonaccessible class.

\begin{claim}{E.2}{Proved}{Enough intrinsic fibrations mean aspherical generation}
For an accessible fundamental localizer $\W$, let $\W^{\mathrm{asph}}$
be the fundamental localizer generated by
$\{A\to e\mid A\text{ is }\W\text{-aspherical}\}$.
The following conditions are equivalent: Loc~(7 F);
$\W=\W^{\mathrm{asph}}$; $\W$ is generated by a small set of
maps $A\to e$; the simplicial localization of \claimref{E.1} is
right proper; and the Thomason model structure on $\Cat$ for
$\W$ is right proper. Thus Loc~(7 F) is a properness condition,
and does not follow merely from existence of a model structure.
\end{claim}
\begin{proof}
Suppose first that (7 F) holds. The class $\Fib_{\W}$ is stable
under composition, pullback, and retracts. All maps to $e$ belong
to it because fundamental localizers are stable under products.
We verify explicitly the only potentially hidden part of the
Brown structure. Let $p:T\to B$ belong to $\Fib_{\W}\cap\W$,
and let $g:X\to B$ be arbitrary. Factor $g$ as
$X\xrightarrow{i}Z\xrightarrow{q}B$, with $i\in\W$ and
$q\in\Fib_{\W}$. Then
$p_Z:T\times_BZ\to Z$ is in $\W$, since $q$ preserves
$\W$ under pullback. The comparison
$T\times_BX\to T\times_BZ$ is also in $\W$, because
$p_Z\in\Fib_{\W}$ and it is the pullback of $i$.
The identity between their composites over $B$, followed by
two out of three, implies $p_X\in\W$. Thus trivial intrinsic
fibrations are universally in $\W$, as required for the Brown
lemma.

For $b\in B$, pull $p$ back along $B/b\to B$. Its total
category is $T/b$. The resulting map belongs to $\W$, and
$B/b$ has a terminal object. Thus $T/b$ is $\W$-aspherical.
By definition its map to $e$ belongs to
$\W^{\mathrm{asph}}$. The Theorem A axiom for
$\W^{\mathrm{asph}}$ now gives $p\in\W^{\mathrm{asph}}$.
The dual Brown lemma, applied to the identity functor with
target weak equivalences $\W^{\mathrm{asph}}$, implies that
every map in $\W$ belongs to $\W^{\mathrm{asph}}$.
The opposite inclusion holds from the definition of the
generators. This direction does not require accessibility.

For accessible $\W$, \citet[Theorem 6.1.11]{RootCisinski2006}
identifies $\W=\W^{\mathrm{asph}}$ with generation by a set
of categories and with properness of the associated simplicial
model structure. \citet[Theorem 5.2.15]{RootCisinski2006} identifies
this with right properness of its Thomason model structure.
In a right proper model structure every model fibration
preserves weak equivalences under strict pullback. Consequently
its factorization of an arbitrary functor as a trivial
cofibration followed by a fibration is a factorization as
$\W$ followed by $\Fib_{\W}$. This proves (7 F). The comparison
with intrinsic fibrations is also stated in
\citet[Example 6.1.14]{RootCisinski2006}.
\end{proof}

\begin{claim}{E.3}{Proved}{The precise modality and local-cartesian dictionary}
For the accessible localization $L$ of \claimref{E.1}, the conditions
of \claimref{E.2} are also equivalent to each of the following:
$L$ is a nullification; its local spaces form a modality; $L$ is
locally cartesian in the sense that
\[
 L(X\times_ZA)\longrightarrow LX\times_ZA                 \tag{E.3.1}
\]
is an equivalence whenever $A,Z$ are local and $X\to Z\leftarrow A$
is a diagram of spaces; and the presentable infinity-category
$\cS_{\W}$ is locally cartesian closed.
Here a modality means the stable orthogonal factorization system
whose right maps have local homotopy fibers and whose left maps
have $L$-aspherical homotopy fibers. Its left class need not equal
the class of maps inverted by $L$.
Nor is its right class identified with the maximal intrinsic
class $\Fib_{\W}$ in $\Cat$: every terminal projection is
an intrinsic fibration, whereas a terminal map is in the
modality's right class exactly when its source is local.
\end{claim}
\begin{proof}
Generation by maps $A_i\to*$ means precisely localization for
which a space $Y$ is local if and only if the constant-map maps
$Y\to\Map(A_i,Y)$ are equivalences. This is nullification.
For completeness its local spaces are closed under dependent
sums over a local base. If $X$ and the spaces $Y_x$ are local,
then for each generator $A_i$ there are natural equivalences
\[
 \Map\!\left(A_i,\int_{x\in X}Y_x\right)
 \simeq\int_{f\in\Map(A_i,X)}\Map_{/X}(A_i,Y)
 \simeq\int_{x\in X}\Map(A_i,Y_x)
 \simeq\int_{x\in X}Y_x.
\]
In the middle expression the map $A_i\to X$ is $f$; after the
constant-map equivalence for $X$, it is constant at $x$. The
equivalences are inverse to constant maps. Thus the reflective
subcategory is closed under dependent sums, which is the
modality condition. This is the argument of
\citet[Theorem 2.19]{CERijkeShulmanSpitters2020}.
The equivalence with the stable factorization-system formulation
is established there in Theorems~1.34 and~1.54, with uniqueness
in Theorem~1.53, and is the definition used
by \citet[Section 3.4]{CEAnelEtAl2020}.

Conversely, an accessible modality presented by a set of maps
$u_i:A_i\to B_i$ is presented by a set of objects: take all
homotopy fibers of the unit maps $A_i\to LA_i$ and
$B_i\to LB_i$. Their indexing spaces are small in the fixed
universe. These fibers are $L$-aspherical. A space $Y$ null for
them is local for both unit maps. Indeed, for either unit
$q:U\to V$ and its fiber diagram $v\mapsto U_v$, the
canonical identification $U\simeq\colim_{v\in V}U_v$ gives
\[
 \Map(U,Y)\simeq\lim_{v\in V}\Map(U_v,Y),\qquad
 \Map(V,Y)\simeq\lim_{v\in V}Y.
\]
The constant-map equivalences $Y\to\Map(U_v,Y)$ are natural
in $v$, including its paths and all higher coherences.
Taking their homotopy limit therefore proves that
$\Map(V,Y)\to\Map(U,Y)$ is an equivalence. This is a
limit over the full base space, with its monodromy, rather
than a product over its set of points. In the commuting square formed
by $u_i$ and its units, $Lu_i$ is an equivalence. Such a space
is therefore $u_i$-local. This proves equality of the local
objects and hence of the localizations. It is exactly
\citet[Lemma 2.20]{CERijkeShulmanSpitters2020}.

For a modality the units have $L$-aspherical fibers. In (E.3.1)
the map $X\times_ZA\to LX\times_ZA$ is a base change of
$\eta_X$ and therefore has $L$-aspherical fibers. It is
inverted by $L$. Its target is local, being a limit of local
objects, so it exhibits the target as the localization of its
source. This proves local cartesianness. Conversely apply
(E.3.1) to the unit $X\to LX$ and a point $*\to LX$.
Every unit fiber is $L$-aspherical. This implies closure of
local objects under dependent sums: for $X\to Z$ with local
$Z$ and all fibers local, (E.3.1) identifies those fibers with
the fibers of $LX\to Z$, so the unit is a fiberwise
equivalence. This is the modality condition. The terminology
and related equivalences agree with
\citet[Propositions 3.1 and 3.3]{CEMacpherson2021}.

Local cartesianness implies local cartesian closedness as
follows. In a slice over a local $Z$, colimits are obtained
by $L$. For local $A\to Z$, (E.3.1), together with universality
of colimits in spaces, shows that pullback to $A$ preserves
colimits in the local slices. The adjoint functor theorem
therefore provides its right adjoint, which is the dependent
product.

The converse uses the special ambient category $\cS$ and
should not be asserted for an arbitrary reflective subcategory
of an arbitrary infinity-topos without qualification.
Every $X\to Z$ in $\cS/Z$ is the colimit of its diagram of
points $*\xrightarrow{x}Z$, indexed by the infinity-groupoid
$X$. These points are local objects of the slice when $Z$ is
local. The slice inclusion $\cS_{\W}/Z\to\cS/Z$ has left
adjoint $L_Z$, taking $X\to Z$ to $LX\to Z$: its adjunction
follows by taking the fibers of the mapping-space adjunction
over the prescribed map to the local space $Z$. Applying
$L_Z$ to the point diagram therefore proves that its colimit
in the local slice is $LX\to Z$. If $\cS_{\W}$ is locally cartesian closed, pullback
along local $A\to Z$ preserves colimits in its slices.
It consequently sends the colimit $LX\to Z$ of that diagram
to the colimit of its fibers $A_x$. The latter are local
because the inclusion of the local category preserves limits.
Computed in spaces, their colimit is $X\times_ZA$.
Localizing this colimit proves (E.3.1). This argument checks
the claimed converse in exactly the setting of the question.
\end{proof}

\begin{claim}{E.4}{Proved}{Perfect categorical models and arbitrary homotopy pullbacks}
For any fundamental localizer $\W$, Loc~(7 $D_{\Winfty}$) is
equivalent to stability of the induced $\W$-equivalences of
spaces under arbitrary homotopy pullback. No accessibility
hypothesis is needed for this equivalence. The test class
$\Parf_{\Winfty}$ is essential: replacing homotopy pullbacks
by arbitrary strict pullbacks of categories would be stronger
and false even for $\Winfty$.
\end{claim}
\begin{proof}
We specify the comparison used in passing between models.
Every functor $v:C\to B$ has a factorization
\[
 C\xrightarrow{i_v}P(v)\xrightarrow{p_v}B,
 \qquad i_v\in\Winfty,\quad p_v\in\Parf_{\Winfty}.
                                                               \tag{E.4.1}
\]
This is the categorical path factorization recalled in XVI,
manuscript pp.~28--29 and 89--90. An independently constructed
factorization with the same properties follows from
\citet[Theorem 5.3.14]{RootCisinski2006}, applied to $\Winfty$:
factor as a trivial cofibration followed by a strongly proper
and smooth Thomason bifibration. We use this latter verified
construction, so do not assume that every detail of the
manuscript's infinite-path presentation is correct.

For $p\in\Parf_{\Winfty}$, every strict base-change square
along $p$ is a homotopy pullback after applying $B$.
To check the scope of this assertion, properness and smoothness
are defined by the fiber-to-comma comparisons in
\citet[Definition 5.3.1]{RootCisinski2006}; both classes are
closed under base change and composition by Corollary~5.3.6
and Proposition~5.3.8 and their duals. Their intersection is
locally constant by Corollary~6.4.8. Pull back $p$ along an
arbitrary $[m]\to B$. For any $k\in[m]$, properness identifies
the strict fiber with the comma category over $k$, and local
constancy identifies this comma category with the one over
$m$. The latter is the total category, because $m$ is terminal
in $[m]$. Thus every point-to-$[m]$ pullback induces a weak
equivalence on the total categories. This is exactly the
finite-ordinal condition in Proposition~6.4.14(iii), which
identifies $p$ as a $\Winfty$-weak fibration. Finally
Proposition~6.2.3 identifies such weak fibrations with the maps
whose every strict base-change square is homotopy cartesian.
This establishes the comparison for every base category,
not only for points.

The dual Brown argument in Part A turns
$\Parf_{\Winfty}\cap\W\subset\W^u$ into
$\Parf_{\Winfty}\subset\Fib_{\W}$, using (E.4.1) and
closure of $\Parf_{\Winfty}$ under composition and base
change. If $u:A\to B$ belongs to $\W$ and $v:C\to B$ is
arbitrary, replace $v$ by (E.4.1). Then
\[
 A\times_BP(v)\longrightarrow P(v)
\]
belongs to $\W$, and its nerve models the homotopy pullback
of $Bu$ along $Bv$. The map $Bi_v$ is the specified comparison
to the original base and is a weak homotopy equivalence.
Minimality and two out of three transfer membership in $\W$
across that comparison.

This suffices for arbitrary diagrams of spaces: present a
space by a simplicial set, and replace it by its category of
simplices. The canonical realization map
$B(\Delta/K)\to|K|$ is a weak equivalence, naturally in
$K$. A homotopy-coherent cospan of spaces may be represented
by a strict cospan of simplicial sets using Kan replacements
and a path object for its homotopies. Applying $\Delta/-$
gives a strict cospan of small categories and a
levelwise weak-equivalent diagram. Its derived pullback is
the specified pullback of spaces, by invariance of the limit
under equivalence of diagrams. Thus the preceding argument
proves homotopy-pullback stability in full scope.

Conversely, if $\W$-equivalences of spaces have that stability,
take $p\in\Parf_{\Winfty}\cap\W$ and any strict base change.
The established nerve comparison identifies it with a
homotopy pullback of $Bp$, hence with a $\W$-equivalence.
This is exactly the universally-$\W$ conclusion in (7 $D_{\Winfty}$).
\end{proof}

\begin{claim}{E.5}{Proved}{The fibration pole has two possibilities}
For a fundamental localizer in the convention of this paper,
Loc~(7 $D_{\Winfty}$) holds if and only if
$\W=\Winfty$ or $\W=\Fl(\Cat)$. Therefore all eight variants
$A,B,C,D,A_V,B_V,C_V,D_V$ of XVI, Section 4.4, for
$\Winfty\subset V\subset\W$, are equivalent in this modern
axiom system and have precisely these two solutions.
Loc~(5) excludes the second solution. This strengthens the
conclusion $\W\subset\Winfty$ of XVI, Theorem 6.1, by
combining it with minimality.
\end{claim}
\begin{proof}
Assume the weakest clause and suppose some $f\in\W$ is not
a weak homotopy equivalence. By \claimref{E.4}, a homotopy
fiber $F$ of $Bf$ is $\W$-aspherical and not weakly
contractible. If it is empty, $\varnothing\to*$ belongs to
$\W$. Otherwise, every chosen point $*\to F$ belongs to
$\W$ by two out of three. Pulling it back along itself shows
that $\Omega F$ is $\W$-aspherical. Iterating this argument,
either $F$ already has more than one component, or a nonzero
homotopy group produces an iterated loop space $Z$ with at
least two components. If there were no such homotopy group
and $F$ were connected, $F$ would be weakly contractible,
contrary to its selection. Choose points $z_0,z_1$ in
distinct components of $Z$. Both $*\to Z$ maps are in
$\W$, while their homotopy pullback is empty. Again
\claimref{E.4} gives $\varnothing\to*\in\W$.

Now $\varnothing\to X$ is in $\W$ for every category $X$:
it is the product of $\varnothing\to e$ with $\id_X$.
For every $u:X\to Y$, the equality
$u\circ(\varnothing\to X)=(\varnothing\to Y)$ and two
out of three give $u\in\W$. Thus $\W$ is total.
If no $f\in\W\setminus\Winfty$ exists, minimality gives
$\W=\Winfty$.

The total class satisfies all the clauses. For $\Winfty$,
Quillen B for a Grothendieck fibration with weakly equivalent
transition functors identifies every base-change square with
a homotopy pullback, so Loc~(7 a) holds. The implications in
Part A now show that both extreme classes satisfy all eight
variants and that any one of those variants implies the
weakest one. Finally $e\amalg e\to e$ proves that the total
class fails Loc~(5).
\end{proof}

\begin{claim}{E.6}{Proved}{Lead 5 with its exact hypotheses}
For an accessible fundamental localizer, the following are
equivalent: Loc~(7); Loc~(7 $D_{\Winfty}$); $L$ preserves
finite limits; and $L$ is either the identity localization of
spaces or the localization to the terminal infinity-category.
These equivalences do not identify Loc~(7 F) with left
exactness. Instead (7 F) has the strictly weaker meanings
in \claimref{E.2} and \claimref{E.3}.
\end{claim}
\begin{proof}
Write $\mathcal E_L$ for the maps of spaces inverted by $L$.
If $L$ preserves finite limits, $\mathcal E_L$ is stable
under homotopy pullback. Conversely suppose it is stable
under homotopy pullback. In a cospan $X\to Z\leftarrow Y$,
factor the canonical map to the pullback of its localized
cospan as
\[
 X\times_ZY\longrightarrow X\times_{LZ}Y
 \longrightarrow LX\times_{LZ}Y
 \longrightarrow LX\times_{LZ}LY.
\]
The last two maps are pullbacks of units, hence are in
$\mathcal E_L$. The first map is the pullback of the diagonal
$Z\to Z\times_{LZ}Z$. This diagonal is in
$\mathcal E_L$, because its projection back to $Z$ is a
pullback of $\eta_Z$ and its composite with the diagonal
is the identity. Thus all three maps are in
$\mathcal E_L$. Their target is local, being a limit of local
objects; it is therefore the localization of the source.
Since $*$ is local, $L$ preserves finite limits. Combining
this calculation with \claimref{E.4} proves the exactness
dictionary, and \claimref{E.5} supplies the classification.

There is also a proof of the classification entirely in
spaces. If a left exact reflector is not the identity, some
unit $X\to LX$ has a noncontractible homotopy fiber. Left
exactness makes that fiber and all its iterated loop spaces
$L$-aspherical. The argument in \claimref{E.5} then shows
$L\varnothing\simeq*$. For any $X$, left exactness applied
to $X\times\varnothing=\varnothing$ gives
$LX\times L\varnothing\simeq L\varnothing$, whence
$LX\simeq*$. This proves that there are exactly two left
exact reflective localizations of spaces; accessibility is
needed here for constructing $L$ from $\W$, not for this
last classification argument.

To see the strict distinction from (7 F), use
$L=\tau_{\leq n}$, $n\geq0$. It is nullification at
$S^{n+1}$ and satisfies (7 F). The homotopy pullback
$*\times_{K(\mathbb Z,n+1)}*\simeq K(\mathbb Z,n)$
is not preserved by $L$, since the localized base is a point.
Thus it is a modality with locally cartesian closed local
category, but it is not left exact. This example also
explains why the reflective equivalences cannot be substituted
for the left class of a general modality. The same distinction
is explicit in \citet[Theorem 3.1 and the discussion after
Corollary 3.9]{CERijkeShulmanSpitters2020} and in the congruence
formalism of \citet{CEAnelEtAl2022}.
\end{proof}

\begin{claim}{E.7}{Proved}{The three classes on XII, manuscript pp.~178--187}
For $n\geq2$ and a map of spaces $f:X\to Y$, use the manuscript's
definitions with their missing range made explicit:
$a_n$ means that every homotopy fiber is nonempty, connected,
and has $\pi_i=0$ for $1\leq i\leq n$;
$b_n$ means that $\tau_{\leq n}f$ is an equivalence;
and $c_n$ means that
\[
 H^i(Y;M)\longrightarrow H^i(X;f^*M)
\]
is an isomorphism for $0\leq i\leq n$, for every local system
of sets at $i=0$, groups at $i=1$, and abelian groups at
$i\geq2$. Then $a_n$ is the left class of the
$n$-truncation modality, whereas $b_n$ is its class of
reflective equivalences. Precisely,
\[
 a_n\subset b_n\subset c_n\subset a_{n-1},
 \qquad a_n=(b_n)^{\mathrm{univ}},
\]
where universal means stable membership under all homotopy
base changes. Moreover $b_2=c_2$, while $b_n\subsetneq c_n$
for every $n\geq3$. Thus the manuscript's expectation that
all displayed implications are strict is false at $n=2$.
The final suggestion on manuscript p.~187 that $a_n$ itself
defines the localizer of $n$-types is also false: $a_n$ fails
two out of three. That localizer is $b_n$.
\end{claim}
\begin{proof}
The long exact sequence of a homotopy fiber gives
\[
 f\in a_n\quad\Longleftrightarrow\quad
 f\in b_n\ \text{and }\pi_{n+1}(X,x)\to\pi_{n+1}(Y,f(x))
 \text{ is surjective at every }x.
\]
This includes the requirement that every component of $Y$
has a fiber. The $n$-connected--$n$-truncated factorization
system factors a map by truncating each homotopy fiber;
the first factor has $n$-connected fibers and the second
has $n$-truncated fibers. Both classes are stable under
base change. The reflective equivalences are the $b_n$,
and their universal subclass is exactly $a_n$: necessity
is the test by points of $Y$, and sufficiency follows from
the unchanged fibers after a base change.

If $f\in b_n$, it induces an equivalence of fundamental
groupoids. For a local system $M$, the twisted
Eilenberg--MacLane space over this fundamental groupoid
representing $H^i(-;M)$ is an $n$-type when $i\leq n$.
Mapping into it therefore depends only on $\tau_{\leq n}$.
This proves $b_n\subset c_n$, including the torsor
interpretation at degree one.

Suppose $f\in c_n$. Degree zero with all constant sets
implies a bijection on components, so we work separately
on each connected component. Degree one with all constant
groups says that
\[
 \Hom(\pi_1Y,G)/G\longrightarrow\Hom(\pi_1X,G)/G
\]
is bijective for every group $G$, with conjugation as the
quotient action. Put $\phi=\pi_1f$. For $G=\pi_1X$,
surjectivity supplies $r:\pi_1Y\to\pi_1X$ with
$r\phi$ conjugate to the identity; conjugate $r$ so that
$r\phi=\id$. For $G=\pi_1Y$, injectivity says that
$\phi r$ is conjugate to $\id$, since both restrict to
$\phi$ on $\pi_1X$. Hence $\phi r$ is an automorphism,
and $\phi$ is surjective as well as injective.

Identify this common group with $\Gamma$. Lift $f$ to an
equivariant map $\widetilde f:\widetilde X\to\widetilde Y$
of universal covers. Choose injective $\mathbb Z[\Gamma]$-modules
$I$. For the free equivariant singular-chain complexes,
injectivity gives
\[
 H^i(X;I)=\Hom_{\mathbb Z[\Gamma]}(H_i(\widetilde X),I),
\]
and the corresponding formula for $Y$. These identifications
are natural in $f$. The degree-$i$ assertions of $c_n$ for
all injectives imply that $H_i(\widetilde f)$ is an
isomorphism for $i\leq n$: an injective cogenerator detects
the kernel and cokernel of a module map. Relative Hurewicz,
applied inductively to the simply connected mapping-cylinder
pair of $\widetilde f$, now gives vanishing relative
homotopy through degree $n$. Equivalently its homotopy
fiber is $(n-1)$-connected. This proves $c_n\subset a_{n-1}$.
At $n=2$, ordinary Hurewicz identifies the second homotopy
groups of the simply connected covers with their second
homology groups. Thus $c_2$ also implies $b_2$.

For strictness, $*\to K(\mathbb Z,n+1)$ is in $b_n$ and
not in $a_n$. Its retraction $K(\mathbb Z,n+1)\to*$ and
their composite $\id_*$ are both in $a_n$, proving the
asserted failure of two out of three. For $n\geq3$, the inclusion
\[
 S^2\vee S^{n-1}\longrightarrow S^2\times S^{n-1}
\]
is in $c_n$: both spaces are simply connected and their
cellular cochain complexes agree through degree $n$ for
every abelian coefficient group. It is not in $b_n$.
The product is obtained by attaching its top $(n+1)$-cell
to the wedge by the Whitehead product
$[\iota_2,\iota_{n-1}]\in\pi_n(S^2\vee S^{n-1})$.
This attaching map is nonzero, for otherwise the product
would be homotopy equivalent to
$S^2\vee S^{n-1}\vee S^{n+1}$, whose positive-dimensional
cup products vanish, whereas the product has a nonzero
cross cup product in degree $n+1$. Its nonzero class is
killed by the inclusion, as required. Finally
$S^n\to*$ belongs to $a_{n-1}$ but fails $c_n$ by degree
$n$ cohomology. These examples also show explicitly why
$c_n$, for $n\geq3$, is neither of the two designated classes
of the $n$-truncation modality; at $n=2$ it is its reflective
equivalence class.
The finite cellular complexes and their integral cup-product
certificates for $3\leq n\leq8$ are checked in
\texttt{code/ce\_truncation\_cochains.py}, with the complete
output in its accompanying JSON file. The code retains the
two distinct degree-two generators when $n=3$; it does not
replace the unstable attaching-map argument by a homology
calculation.
\end{proof}

The local coefficients in this last proof may be described
covariantly as left $\mathbb Z[\Gamma]$-modules or, with
the manuscript's presheaf convention, as right modules
using the inverse action. The equivariant-chain and
cochain identifications use one of these choices consistently;
no change in variance is being hidden in the comparison.
